\section{集中解析}

\needspace{10\baselineskip}
\begin{solution}{题 1 · 解析}
\think 用旋转连接半角与线段和，用垂线连接等角与等距。
\diagram{\figTwoEquivalence}

\solve 将 \(M\) 绕 \(B\) 顺时针旋转 \(90^\circ\) 到 \(P\)，得 \(CP=AM\)、\(BP=BM\)、\(\angle MBP=90^\circ\)；另作 \(BQ\perp MN\)。两种辅助线分图表示。
\diagram{\figTwoSolve}

\solvehead{旋转：半角与线段和相互推出}
\((1)\) 使 \(BN\) 平分 \(\angle MBP\)，因此
\[
\triangle MBN\cong\triangle PBN\ (\mathrm{SAS})
\ \Longrightarrow\ \relax
\begin{cases}
MN=PN=AM+CN &(2),\\
\angle BMN=\angle BPN=\angle AMB &(3),\\
\angle BNM=\angle BNP=\angle BNC &(4).
\end{cases}
\]
反过来，\((2)\) 给出 \(MN=PN\)，结合 \(BM=BP\)、\(BN\) 公共，
\[
\triangle MBN\cong\triangle PBN\ (\mathrm{SSS})
\quad\Longrightarrow\quad
\angle MBN=\tfrac12\angle MBP=45^\circ\quad(1).
\]

\solvehead{垂线：两个等角条件相互推出}
\(B\) 在 \(\angle AMN\)、\(\angle CNM\) 内，角平分线可用到两边的等距来判断。
\[
\begin{aligned}
(3)&\ \Longrightarrow\ \mathrm{Rt}\triangle BQM\cong\mathrm{Rt}\triangle BAM
\ \Longrightarrow\ BQ=BA=BC\ \Longrightarrow\ (4);\\
(4)&\ \Longrightarrow\ BQ=BC=BA
\ \Longrightarrow\ BM\text{ 平分 }\angle AMN\ \Longrightarrow\ (3).
\end{aligned}
\]
第一行的全等使用公共斜边 \(BM\) 与 \(M\) 处相等的锐角。

\solvehead{等角：恢复顶点处的半角}
令 \(\alpha=\angle ABM\)、\(\beta=\angle NBC\)、\(\theta=\angle MBN\)。由 \((3)(4)\)，
\[
\left.\begin{aligned}
\theta&=\alpha+\beta &&(\triangle BMN),\\
\theta+\alpha+\beta&=90^\circ &&(\angle ABC)
\end{aligned}\right\}
\quad\Longrightarrow\quad \theta=45^\circ\quad(1).
\]
\insight{旋转看全等，垂线看等距；最后用三角形内角和恢复半角。}
\end{solution}
\needspace{10\baselineskip}
\begin{solution}{题 2 · 解析}
\think 题面没有写 \(45^\circ \)，但已给出题 \(1\) 的等价条件。

\solve 因 \(AD\parallel BC\)，\(\angle AMB=\angle MBC\)。已知 \(\angle BMN=\angle MBC\)，故 \(\angle AMB=\angle BMN\)，由题 \(1\)，\(\angle MBN=45^\circ \)且 \(MN=AM+CN\)。\(AM=9-6=3\)。设 \(CN=x\)，则 \(DN=9-x\)，\(MN=3+x\)。在 \(\mathrm{Rt}\triangle MDN\) 中，\(6^2+(9-x)^2=(3+x)^2\)，解得 \(x=\frac{9}{2}\)。所以 \(CN=\frac{9}{2}\)。

\insight{没有出现半角数值时，先找模型的等价条件。}
\end{solution}
\needspace{10\baselineskip}
\begin{solution}{题 3 · 解析}
\think \(E\) 在正方形对角线上，天然具有对称性；先把 \(BE\) 变成 \(DE\)，再用垂直关系证明 \(DE=EN\)。

\diagram{\figFourSolve}
\solve 连接 \(DE\)。正方形关于 \(AC\) 对称，\(B\) 与 \(D\) 互为对称点，\(E\) 位于对称轴上，所以 \(BE=DE\)、\(\angle ABE=\angle ADE\)。由 \(BE\perp EN\)、\(BC\perp ND\)，且按图中射线方向，有 \(\angle DNE=\angle EBC\)。又 \(\angle EDN=90^\circ -\angle ADE=90^\circ -\angle ABE=\angle EBC\)，所以 \(\angle EDN=\angle DNE\)，\(DE=EN\)。因此 \(BE=EN\)。

\insight{对角线给出等长，垂直给出等角，合起来构成等腰三角形。}
\end{solution}
\needspace{10\baselineskip}
\begin{solution}{题 5 · 解析}
\think 长边是短边的两倍，可以截出一个正方形，先在小图中求长，再回到原矩形。

\diagram{\figSixSolve}
\solve 取 \(AB\) 中点 \(M\)，作 \(MQ\parallel AD\)，交 \(DC\) 于 \(Q\)，\(AE\) 与 \(MQ\) 交于 \(N\)。\(AMQD\) 是边长 \(6\) 的正方形，\(DQ=6\)、\(FQ=3\)。设 \(MN=x\)，图 A 给出 \(FN=DF+MN=3+x\)，而 \(NQ=6-x\)。在 \(\mathrm{Rt}\triangle FQN\) 中，\((3+x)^2=3^2+(6-x)^2\)，得 \(x=2\)。\(\triangle AMN\sim \triangle ABE\)，\(\frac{AB}{AM}=2\)，因此 \(BE=2MN=4\)。

\insight{矩形先截正方形，局部模型求出后再用相似放回整体。}
\end{solution}
\needspace{16\baselineskip}
\begin{solution}{题 6 · 解析}
\think 交点处的 \(45^\circ \)可通过一条垂线搬到 \(A\)，随后变成例题 \(4\) 的底边高模型。

\diagram{\figSevenSolve}
\solve 过 \(A\) 作 \(AG\perp BD\)，交 \(BC\) 向 \(C\) 外的延长线于 \(G\)。\(\angle BFE=45^\circ \)，所以按 \(G\)、\(C\)、\(E\) 的点序，\(\angle GAE=45^\circ \)。\(\mathrm{Rt}\triangle ACG\) 与 \(\mathrm{Rt}\triangle BCD\) 相似，\(\frac{CG}{AC}=\frac{CD}{BC}=\frac{3}{8}\)，故 \(CG=\frac{9}{4}\)。设 \(CE=x\)。在 \(\triangle AGE\) 中，\(AC\) 为高，\(AC=6\)、\(CG=\frac{9}{4}\)。按例题 \(4\) 的延长旋转构造，\((\frac{9}{4}+x)^2=(6-\frac{9}{4})^2+(6-x)^2\)，展开并合并同类项，得 \(x=\frac{30}{11}\)。所以 \(CE=\frac{30}{11}\)。

\insight{交叉角优先用平行或垂直转移，再进入熟悉的模型。}
\end{solution}
\clearpage
\begin{solution}{题 7 · 解析}
\think 先用旋转得到相等的斜边，再从图上读出“一差一和”，用勾股定理求值。

\solve 设 \(CE=CF=x\)，则 \(AB=x+3\)、\(AD=x+4\)。将 \(E\) 绕 \(A\) 逆时针旋转 \(90^\circ\) 到 \(G\)。作 \(GH\parallel AD\)，交 \(CD\) 的延长线于 \(H\)；作 \(GT\perp AD\)，垂足为 \(T\)。
\diagram{\figEightSolve}

\solvehead{旋转全等：两条斜边相等}
\[
\left.\begin{gathered}
AE=AG,\quad AF\text{ 公共}\\
\angle EAF=\angle GAF=45^\circ
\end{gathered}\right\}
\ \xRightarrow{\mathrm{SAS}}\ \triangle EAF\cong\triangle GAF
\ \Longrightarrow\ EF=GF.
\]

\solvehead{读图：横向取差，竖向取和}
旋转给出 \(AT=AB=x+3\)、\(GT=BE=4\)；\(GTDH\) 为矩形，故
\[
\begin{aligned}
GH&=TD=AD-AT=(x+4)-(x+3)=1,\\
FH&=FD+DH=3+4=7.
\end{aligned}
\]

\solvehead{勾股：同一斜边的平方相等}
\[
\left.\begin{aligned}
EF^2&=CE^2+CF^2=2x^2,\\
GF^2&=GH^2+FH^2=1^2+7^2=50
\end{aligned}\right\}
\ \xRightarrow{EF=GF}\ 2x^2=50.
\]
\[
x>0\quad\Longrightarrow\quad x=5
\quad\Longrightarrow\quad \boxed{AB=x+3=8}.
\]
\insight{相等的 \(CE\)、\(CF\) 让横向的 \(x\) 抵消，竖向直接相加。}
\end{solution}
\needspace{16\baselineskip}
\begin{solution}{题 8 · 解析}
\think 辅助线完全相同，位置改变只影响共线线段的加减。

\diagram{\figNineSolve}
\solve \((1)\)将 \(M\) 绕 \(A\) 逆时针旋转 \(90^\circ \)到 \(P\)，则 \(P\) 在 \(CD\) 向 \(D\) 外的延长线上，\(DP=BM\)、\(AP=AM\)。\(\angle MAP=90^\circ \)，\(AN\) 平分此角，所以 \(\triangle MAN\cong \triangle PAN(\mathrm{SAS})\)，\(MN=PN\)。由 \(P\)、\(D\)、\(N\) 的点序，\(PN=PD+DN=BM+DN\)。

\((2)\)同样旋转 \(M\) 得到 \(P\)。此时 \(P\)、\(N\)、\(D\)、\(C\) 按此顺序共线。\(\angle MAP=90^\circ \)、\(\angle MAN=45^\circ \)，\(AN\) 仍在该直角内部并平分它，所以 \(\triangle MAN\cong \triangle PAN\)，\(MN=PN\)。由点序，\(PN=PD-ND=BM-DN\)，即 \(BM=MN+DN\)。这里 \(BM>DN\)，不能把两条正长度的和照搬到外部构型。

\insight{先证明全等，再根据点序写加法或减法。}
\end{solution}
\clearpage
\begin{solution}{题 9 · 解析}
\think 旋转后得到一个 \(3\!:\!4\!:\!5\) 的直角三角形，再按 \(K\) 在 \(C\) 的哪一侧分两种情况。

\solve 作 \(AH\perp BC\)。\(\angle B=45^\circ\) 使 \(H\) 在 \(B\) 向 \(C\) 的射线上，设 \(AH=BH=t>0\)。将 \(\triangle ABD\) 绕 \(A\) 逆时针旋转 \(90^\circ\) 到 \(\triangle AKG\)，\(B,D\) 分别对应 \(K,G\)。
\diagram{\figTenSolve}

\solvehead{旋转：搬出已知的直角边}
\[
\begin{aligned}
\triangle ABH\text{ 为等腰直角三角形}
&\ \Longrightarrow\ BH=HK=t\ \Longrightarrow\ BK=2t;\\
BD\xrightarrow[\text{绕 }A]{90^\circ}KG
&\ \Longrightarrow\ KG\perp BC,\quad KG=BD=4.
\end{aligned}
\]

\solvehead{半角全等：搬出斜边}
\[
\left.\begin{gathered}
AD=AG,\quad AC\text{ 公共}\\
\angle DAC=\angle CAG=45^\circ
\end{gathered}\right\}
\ \xRightarrow{\mathrm{SAS}}\ \triangle DAC\cong\triangle GAC
\ \Longrightarrow\ GC=DC=5.
\]

\solvehead{勾股与分类：先求一段，再分加减}
\[
\mathrm{Rt}\triangle GKC\quad\Longrightarrow\quad
CK^2=GC^2-KG^2=5^2-4^2=9\quad\Longrightarrow\quad CK=3.
\]
由 \(BK=2t\)、\(AB=\sqrt2\,t\)，两种位置分别给出：
\[
\begin{array}{c@{\qquad\qquad}c}
K\text{ 在 }C\text{ 左侧}&K\text{ 在 }C\text{ 右侧}\\
\Downarrow&\Downarrow\\
BK=9-3=6&BK=9+3=12\\
\Downarrow&\Downarrow\\
t=3&t=6\\
\Downarrow&\Downarrow\\
\boxed{AB=3\sqrt2}&\boxed{AB=6\sqrt2}
\end{array}
\]
两种构型都满足原条件；若限定 \(H\) 在 \(DC\) 内部，才只保留 \(6\sqrt2\)。
\insight{先用旋转固定 \(CK=3\)，再由点序决定 \(BK\) 用加法还是减法。}
\end{solution}
